\chapter{Convex bodies}\label{geometry:ch:convex}

\begin{center}\itshape
They bought as much land as they could enclose with the hide of a bull.
\upshape(paraphrase of Virgil, \emph{Aeneid}, Book~I, lines 367--368)
\end{center}

Dido's bargain in the \emph{Aeneid} turns a purchase of land into a problem about length and area. The legend has her cut a bull's hide into narrow strips and join them into a cord. Given the cord, which boundary encloses the most land? A circle is a persuasive guess, but a picture does not prove that every competing shape encloses less.

Volume~5, Theorem~6.17 gave Hurwitz's Fourier-series proof of the planar isoperimetric inequality. Here another operation on shapes leads to an inequality in every positive dimension. To add two sets, take all sums of one point from each. The Brunn--Minkowski inequality controls the volume of that sum. Adding a small ball then relates the increase in volume to surface area.

Convexity makes this operation easier to understand. A convex set contains every segment between its points. A disk and a triangle are convex; a crescent is not. We shall use that simple segment condition to decide when many sets share a point, measure perimeters by shadows, and compare shapes of constant width. The measure and probability prerequisites are Lebesgue measure from Volume~7 and elementary probability from Volume~4. Chapter~11 has developed the general curve versions of the Buffon and Cauchy formulas; the polygon arguments below can also be read on their own.

\begin{think}
Drop a straight needle of length $3/4$ on parallel lines spaced $1$ apart, with uniform position and direction. Is crossing more likely than missing? If you bend the needle, does its crossing probability still depend only on length, or is it a different quantity that remains unchanged?
\end{think}

\section{Convex sets and Helly's theorem}

A set $C\subset\R^n$ is \emph{convex} when $(1-t)x+ty\in C$ for all $x,y\in C$ and $0\le t\le1$. The empty set is convex, and intersections of convex sets are convex. The \emph{convex hull} $\operatorname{conv}(S)$ is the smallest convex set containing $S$. It consists exactly of finite convex combinations $\sum_i t_ix_i$, where $x_i\in S$, $t_i\ge0$ and $\sum_i t_i=1$. These combinations form a convex set and belong to every convex set containing $S$, which proves the description. The hull of the empty set is empty.

In the plane, three points always have an affine dependence once a fourth is added. Separating the positive and negative coefficients turns such a dependence into a point shared by two hulls.

\begin{lemma}[Radon, 1921]\label{geometry:lem:radon}
For $m\ge n+2$ points $x_1,\ldots,x_m\in\R^n$, the index set can be partitioned into two nonempty sets $I,J$ whose corresponding convex hulls meet.
\end{lemma}

\begin{proof}
There are more than $n+1$ vectors $(x_i,1)\in\R^{n+1}$, so choose a nonzero relation
$\sum_i\lambda_ix_i=0$, $\sum_i\lambda_i=0$. Some coefficients are positive and some negative. Put their positive indices in $I$ and all remaining indices, including any zero coefficients, in $J$. The number
$s=\sum_{i\in I}\lambda_i=\sum_{j\in J}(-\lambda_j)$ is positive. Dividing the relation by $s$ gives
\[
\sum_{i\in I}\frac{\lambda_i}{s}x_i
=\sum_{j\in J}\frac{-\lambda_j}{s}x_j.
\]
Each side is a convex combination, so their common value lies in both hulls. Repeated points cause no difficulty: the partition is of indices.
\end{proof}

\begin{theorem}[Helly]\label{geometry:thm:helly}
Let $C_1,\ldots,C_m\subset\R^n$ be convex, with $m\ge n+1$. If every $n+1$ of these sets have a common point, then all $m$ have a common point.
\end{theorem}

\begin{proof}
Induct on $m$. At $m=n+1$ the hypothesis is the conclusion. For $m\ge n+2$, induction says that every subfamily of $m-1$ sets intersects, because each of its $n+1$ subfamilies does. Choose $x_i\in\bigcap_{j\ne i}C_j$ for each omitted index $i$.

Apply Radon's lemma to these points and let $p$ belong to both resulting hulls. Fix $k$. One of the two index sets does not contain $k$. All the points indexed by that set lie in $C_k$, so their convex hull lies in $C_k$ and contains $p$. This works for every $k$, placing $p$ in the full intersection.
\end{proof}

No closedness or boundedness was used; the family must be finite. For $n=0$ the hypothesis merely says that every set contains the unique point of $\R^0$, so the conclusion still holds. Exercise~\ref{geometry:ex:hellysharp} explains why $n+1$ cannot be reduced for positive dimension.

Helly found the result in 1913 and published it in 1923, after years as a prisoner of war in Siberia. Radon's proof appeared in 1921. The historical account by Danzer, Gr\"unbaum and Klee is listed below. The useful feature of the theorem is that a question about an entire finite family reduces to small subfamilies, independently of how many sets there are.

\section{Buffon's needle}

Put horizontal lines at every integer height. To throw a fixed polygonal curve, choose a rotation angle $\theta$ uniformly in $[0,2\pi)$ and a vertical translation $u$ uniformly in $[0,1)$, independently. Horizontal translation cannot affect crossings. Let $N$ count intersections with the ruled lines, with traversal multiplicity if an arc is retraced. A crossing at a vertex can be counted once; positions putting a vertex on a ruled line have probability zero. A segment lying along a ruled line is another null event, so its exceptional intersection count can be assigned arbitrarily.

Buffon posed geometric throwing problems to the Paris Academy in 1733 and published his needle calculation in 1777. The model requires explicit choices of position and direction: without them, a phrase such as ``at random'' does not specify a probability.

\begin{lemma}\label{geometry:lem:randominterval}
For $h\ge0$ and uniform $u\in[0,1)$, the number of integers in $[u,u+h)$ has expectation $h$.
\end{lemma}

\begin{proof}
Write the count as $\sum_{k\in\Z}\mathbf1_{[u,u+h)}(k)$. Integrate in $u$ and translate each term using $v=k-u$. Up to endpoints of measure zero, the intervals $(k-1,k]$ partition the real line. Hence
\[
\int_0^1\sum_{k\in\Z}\mathbf1_{[0,h)}(k-u)\,du
=\int_{\R}\mathbf1_{[0,h)}(v)\,dv=h.
\]
Tonelli applies because the summands are nonnegative. The case $h=0$ gives an empty interval and count zero.
\end{proof}

\begin{theorem}[Buffon, Barbier]\label{geometry:thm:buffon}
A polygonal curve of total length $L$, thrown by the model above, satisfies $\mathbb E[N]=2L/\pi$. In particular, a straight needle of length $0\le\ell\le1$ crosses a line with probability $2\ell/\pi$.
\end{theorem}

\begin{proof}
Count crossings segment by segment and add their expectations. A segment of length $\ell$ with direction angle $\theta$ has vertical extent $\ell|\sin\theta|$. Conditional on this angle, the interval lemma gives the same number as its mean crossing count. An offset of the segment's endpoint only translates the uniform position modulo $1$, leaving the expectation unchanged. Average over angles:
\[
\mathbb E[N]=\frac{\ell}{2\pi}\int_0^{2\pi}|\sin\theta|\,d\theta
=\frac{2\ell}{\pi}.
\]
Summing proves the curve formula. For a straight needle of length at most $1$, two distinct line crossings would require vertical extent at least $1$, an equality event of probability zero at the endpoint $\ell=1$. Thus $N$ is $0$ or $1$ almost surely and its expectation is its crossing probability. At $\ell=0$ that probability is zero.
\end{proof}

A bent needle of fixed length has the same expected number of crossings, as Barbier observed. Its probability of at least one crossing need not be the same: some placements can create several crossings. For the opening needle, $2\ell/\pi=3/(2\pi)<1/2$, so missing is slightly more likely. A circle of diameter $1$ meets the lines twice for almost every vertical placement; its length $\pi$ is consistent with the general curve formula of Chapter~11.

\begin{figure}[ht]
\centering
\fig{geometry-ch11-buffon}
\caption{Forty needles of length $0.8$ thrown on lines at distance $1$. Here $20$ of them cross a line (in rust); the theorem predicts a proportion
$\frac{1.6}\pi\approx0.51$.}
\label{geometry:fig:buffon}
\end{figure}

For $T$ independent throws and $C>0$ crossings of a short straight needle, the usual estimate is $\pi\approx2\ell T/C$. Lazzarini reported $3408$ throws with $\ell=5/6$ and the estimate $355/113$ in 1901. Badger's 1994 analysis explains why these numbers invite suspicion; Exercise~\ref{geometry:ex:lazzarini} exposes their arithmetic. Agreement to many digits is not by itself evidence of a sound experiment. The sampling error normally shrinks on the scale $T^{-1/2}$, so reducing typical error tenfold requires about one hundred times as many throws.

\section{Width and curves of constant width}

For a nonempty compact convex set $K$ in the plane, its \emph{width} $w_K(\theta)$ is the length of its orthogonal projection onto the line with direction $\theta$. Equivalently, it is the distance between the two parallel support lines perpendicular to that direction. A \emph{support line} touches the set while leaving it wholly on one side.

For a set with interior, define its perimeter $L(K)$ as the supremum of the perimeters of convex polygons contained in it. This agrees with boundary length. For a segment of length $a$ use $L=2a$, counting the trip along it and back; for a point use $L=0$. These conventions make the formulas continuous as a thin polygon collapses. Empty sets have no width here.

\begin{theorem}[Cauchy's formula]\label{geometry:thm:cauchy}
Every nonempty compact convex planar set satisfies
\[
L(K)=\int_0^\pi w_K(\theta)\,d\theta.
\]
Thus perimeter is $\pi$ times mean width.
\end{theorem}

\begin{proof}
For a convex polygon with edge lengths $\ell_i$ and edge directions $\theta_i$, projecting its boundary covers the projection interval twice, once in each direction. Consequently
\[
w_K(\theta)=\frac12\sum_i\ell_i|\cos(\theta-\theta_i)|.
\]
The integral of each absolute cosine over an interval of length $\pi$ is $2$, giving the formula for polygons.

For a general $K$ with interior, take nested inscribed polygons $P_j$ as follows. Include increasingly fine finite subsets of a dense subset of $K$, and also the vertices of polygons whose perimeters approach the defining supremum. Taking their convex hulls ensures $L(P_j)\to L(K)$. This supremum is finite, since the polygon formula bounds every perimeter by $\pi\operatorname{diam}K$. Compactness and density give numbers $\delta_j\to0$ with $K\subset P_j+\delta_j\mathbb B$. Projection then yields
\[
0\le w_K(\theta)-w_{P_j}(\theta)\le2\delta_j
\]
uniformly in $\theta$. Integrating and taking limits proves the formula. For a segment, $w_K(\theta)=a|\cos(\theta-\theta_0)|$, whose integral is $2a$; for a point it is zero.
\end{proof}

A set has \emph{constant width} $w$ if its width equals $w$ in every direction. A disk of diameter $w$ does. So does a Reuleaux triangle: intersect the three radius-$w$ disks centred at the vertices of an equilateral triangle of side $w$. Its boundary consists of three circular arcs. In a support direction facing an arc, the opposite support is at that arc's centre vertex, a distance $w$ away; the remaining directions exchange the roles or give the common endpoints of the arcs. Rotating the direction covers all three arc--vertex pairs.

\begin{corollary}[Barbier, 1860]\label{geometry:cor:barbier}
A nonempty compact convex set of constant width $w$ has perimeter $\pi w$.
\end{corollary}

\begin{proof}
Cauchy's formula integrates the constant function $w$ over $[0,\pi]$.
\end{proof}

\begin{figure}[ht]
\centering
\fig{geometry-ch11-reuleaux}
\caption{A Reuleaux triangle and a disk of constant width $1$, each turning inside a unit square while touching its four sides. Both have perimeter $\pi$.}
\label{geometry:fig:reuleaux}
\end{figure}

The equal width keeps a platform resting on such rollers at a constant height. It does not say that a fixed axle at the roller's centre would stay at constant height. Reuleaux studied mechanisms of this kind in the nineteenth century; related mechanisms allow a rotating Reuleaux triangle to trace an almost square hole. We quote the Blaschke--Lebesgue theorem that the Reuleaux triangle minimizes area at fixed positive width. The opposite area bound, attained by the disk, follows in Exercise~\ref{geometry:ex:constwidthdisk}.

\section{The Brunn--Minkowski inequality}

For sets $A,B\subset\R^n$, their \emph{Minkowski sum} is
\[
A+B=\{a+b:a\in A,\ b\in B\}.
\]
Two intervals add by adding their endpoints. Adding the unit ball $\mathbb B$ scaled by $\varepsilon\ge0$ to a nonempty compact set gives all points at distance at most $\varepsilon$ from it. A closest point exists by compactness, which proves this last description. Sums of compact sets are compact because addition is continuous; sums of convex sets are convex by distributing convex combinations.

\begin{figure}[ht]
\centering
\fig{geometry-ch11-minkowski}
\caption{The Minkowski sum of a triangle $A$ (area $1.4$) and a square $B$ (area $0.49$) is a hexagon of area $4.27$, and
$\sqrt{4.27}\approx2.07\ge\sqrt{1.4}+\sqrt{0.49}\approx1.88$.}
\label{geometry:fig:minkowskisum}
\end{figure}

\begin{figure}[ht]
\centering\fig{v9-app-j-helly-sharp}
\caption{The three edges of a triangle meet in pairs, but all three have no common point. Convexity alone cannot lower Helly's planar requirement from three sets to two.}
\label{fig:v9-app-j-helly-sharp}
\end{figure}

\begin{figure}[ht]
\centering\fig{v9-app-j-parallel-body}
\caption{Adding a radius-$\varepsilon$ disk to a rectangle adds four side strips and four quarter-disks. Their areas give $A+L\varepsilon+\pi\varepsilon^2$; the same decomposition uses exterior-angle sectors for a general polygon.}
\label{fig:v9-app-j-parallel-body}
\end{figure}

\begin{theorem}[Brunn--Minkowski]\label{geometry:thm:bm}
For $n\ge1$ and nonempty compact $A,B\subset\R^n$,
\[
\operatorname{vol}(A+B)^{1/n}
\ge\operatorname{vol}(A)^{1/n}+\operatorname{vol}(B)^{1/n}.
\]
No convexity or positive-volume assumption on $A,B$ is required.
\end{theorem}

\begin{proof}
Write $|X|$ for volume. We first treat finite unions of closed axis-parallel boxes with positive side lengths and pairwise disjoint interiors, inducting on the total number of boxes in the two unions.

For two boxes, let the side lengths be $a_i,b_i>0$. Their sum has sides $a_i+b_i$. Arithmetic--geometric mean gives
\[
\prod_{i=1}^n\left(\frac{a_i}{a_i+b_i}\right)^{1/n}
+\prod_{i=1}^n\left(\frac{b_i}{a_i+b_i}\right)^{1/n}
\le\frac1n\sum_{i=1}^n\frac{a_i+b_i}{a_i+b_i}=1.
\]
Multiplying by $|A+B|^{1/n}$ proves the base case.

If at least three boxes occur, suppose $A$ has at least two. Two of its boxes have disjoint interiors, so in some coordinate their open intervals are disjoint. Choose a coordinate hyperplane $x_j=t$ between those intervals. Cut all boxes along it and discard pieces with zero volume. Each resulting side $A^+,A^-$ contains positive volume and has fewer boxes than $A$: one of the two selected boxes is entirely excluded on each side. Each original box contributes at most one piece to either side.

Put $\lambda=|A^+|/|A|\in(0,1)$. Choose $s$ so that cutting $B$ at $x_j=s$ puts the same volume fraction $\lambda$ on its positive side. Such an $s$ exists because the upper-side volume is a continuous function ranging from $|B|$ to $0$. The resulting $B^+,B^-$ have at most as many boxes each as $B$. The sums $A^++B^+$ and $A^-+B^-$ lie on opposite sides of $x_j=t+s$, overlap only in a null hyperplane, and are subsets of $A+B$. Discarded zero-volume cut pieces are not needed for this lower bound. Induction applies to each pair and gives
\[
\begin{aligned}
|A+B|&\ge|A^++B^+|+|A^-+B^-|\\
&\ge\lambda\bigl(|A|^{1/n}+|B|^{1/n}\bigr)^n
+(1-\lambda)\bigl(|A|^{1/n}+|B|^{1/n}\bigr)^n.
\end{aligned}
\]
This proves the claim for box unions.

For compact $A,B$, take the unions $A_j,B_j$ of closed dyadic grid cubes of side $2^{-j}$ meeting the respective sets. There are finitely many cubes; their interiors are disjoint. The unions decrease to $A,B$, so continuity of finite measure from above gives $|A_j|\to|A|$ and $|B_j|\to|B|$. Every point of $A_j$ is within $\delta_j=\sqrt n\,2^{-j}$ of $A$, and likewise for $B_j$. Hence
\[
A_j+B_j\subset(A+B)+2\delta_j\mathbb B.
\]
The sets on the right decrease to the compact set $A+B$. Another use of continuity from above now gives
\[
|A+B|\ge\limsup_j|A_j+B_j|
\ge\lim_j\bigl(|A_j|^{1/n}+|B_j|^{1/n}\bigr)^n
=\bigl(|A|^{1/n}+|B|^{1/n}\bigr)^n.
\]
This passage also covers zero-volume sets.
\end{proof}

Brunn established his convex-body inequality in 1887, and Minkowski developed the theory further. The box-cutting proof used here is due to Hadwiger and Ohmann in 1956; Gardner's survey traces these versions. The matching volume fractions are the reason the induction works: both contributions become the same expression multiplied by complementary fractions.

\begin{corollary}[isoperimetric inequality]\label{geometry:cor:isoperimetric}
For $n\ge1$ and a compact set $K\subset\R^n$ of positive volume, define its outer surface content by
\[
S(K)=\liminf_{\varepsilon\to0^+}
\frac{|K+\varepsilon\mathbb B|-|K|}{\varepsilon}.
\]
Then $S(K)\ge n|\mathbb B|^{1/n}|K|^{(n-1)/n}$, and balls attain equality. For a planar compact convex set of area $A$ and perimeter $L$, this gives $L^2\ge4\pi A$.
\end{corollary}

\begin{proof}
Apply Brunn--Minkowski to $K$ and $\varepsilon\mathbb B$:
\[
|K+\varepsilon\mathbb B|
\ge\bigl(|K|^{1/n}+\varepsilon|\mathbb B|^{1/n}\bigr)^n
=|K|+n\varepsilon|K|^{(n-1)/n}|\mathbb B|^{1/n}+O(\varepsilon^2).
\]
Subtract $|K|$, divide by $\varepsilon>0$, and take the lower limit. For $K=r\mathbb B$, $r>0$, direct calculation gives $S(K)=nr^{n-1}|\mathbb B|$, exactly the lower bound. In the plane, $|\mathbb B|=\pi$. Exercise~\ref{geometry:ex:steiner} proves $S(K)=L(K)$ for a compact convex set with interior; its conclusion turns $S\ge2\sqrt{\pi A}$ into the perimeter inequality. For a segment or point $A=0$, that inequality is immediate from the perimeter convention.
\end{proof}

Thus a ball is a minimizer of surface content at fixed positive volume. The proof establishes an inequality for every compact set without first assuming that an optimal set exists. It does not classify every equality case for this broad surface-content definition. In the planar convex problem, the disk attains the bound and therefore encloses the maximum area for its perimeter. This is the part of Dido's question supplied by the inequality; no numerical experiment can replace that comparison with every competitor.

\section{Exercises}

\begin{exercise}[$\star$]\label{geometry:ex:helly1d}
Prove the one-dimensional Helly theorem directly: a finite nonempty family of nonempty intervals meeting pairwise has a common point. Allow unbounded, open and half-open intervals as well as closed ones.
\end{exercise}

\begin{exercise}[$\star$]\label{geometry:ex:hellysharp}
For each $n\ge0$, construct $n+1$ convex sets in $\R^n$ such that every $n$ meet but the full family does not.
\end{exercise}

\begin{exercise}[$\star\star$]\label{geometry:ex:hellydisk}
Let $S$ be a finite set of planar points. Suppose that every subset of at most three points fits in a closed disk of radius $1$, with the centre allowed to depend on the subset. Prove that one radius-$1$ disk contains all of $S$. Explain why saying only ``every three points'' would fail for a two-point set.
\end{exercise}

\begin{exercise}[$\star\star$]\label{geometry:ex:caratheodory}
Prove Carath\'eodory's theorem: for $S\subset\R^n$, every point of $\operatorname{conv}(S)$ is a convex combination of at most $n+1$ points of $S$. Include $n=0$ and explain the empty-set case.
\end{exercise}

\begin{exercise}[$\star$]\label{geometry:ex:lazzarini}
For $3408$ throws of a straight needle of length $5/6$ on unit-spaced lines, how many crossings produce the estimate $355/113$? Find the smallest positive throw count that could produce that exact fraction. What does this say, and what does it not prove, about the reported experiment?
\end{exercise}

\begin{exercise}[$\star$]\label{geometry:ex:noodle}
Find the expected crossing count for the boundary of a square of side $a>0$ and a circle of radius $r>0$ on unit-spaced lines. For which $a$ does the square meet a line in every placement and direction? Justify the circle calculation directly, or use the curve formula already proved in Chapter~11.
\end{exercise}

\begin{exercise}[$\star$]\label{geometry:ex:reuleauxarea}
Find the area of a width-$1$ Reuleaux triangle. Compare it with a disk of the same width, and state the scaled formula for width $w>0$.
\end{exercise}

\begin{exercise}[$\star\star$]\label{geometry:ex:minkowskiwidth}
For nonempty compact convex planar sets $K,M$, prove $w_{K+M}=w_K+w_M$ and $L(K+M)=L(K)+L(M)$. Deduce the perimeter of $K+\varepsilon\mathbb B$ for $\varepsilon\ge0$, including when $K$ is a segment or point.
\end{exercise}

\begin{exercise}[$\star\star$]\label{geometry:ex:steiner}
Prove Steiner's formula $|K+\varepsilon\mathbb B|=A+L\varepsilon+\pi\varepsilon^2$ for a convex polygon of area $A$ and perimeter $L$, for all $\varepsilon\ge0$. Extend it rigorously to every nonempty compact convex planar set. Conclude that the surface content in Corollary~\ref{geometry:cor:isoperimetric} equals perimeter when $K$ has interior.
\end{exercise}

\begin{exercise}[$\star\star$]\label{geometry:ex:constwidthdisk}
Show that a nonempty compact convex set of constant width $w$ has area at most $\pi w^2/4$, and that a disk of diameter $w$ attains this bound. Do not assume an equality-case classification that has not been proved.
\end{exercise}

\begin{exercise}[$\star$]\label{geometry:ex:bmone}
Prove directly that $|A+B|\ge|A|+|B|$ for nonempty compact subsets $A,B$ of the real line, including sets of measure zero.
\end{exercise}

\begin{exercise}[$\star\star$]\label{geometry:ex:bmconcave}
For a convex set $A$ and $0\le\lambda\le1$, prove $(1-\lambda)A+\lambda A=A$. For $n\ge1$ and nonempty compact $A,B\subset\R^n$, deduce
\[
|(1-\lambda)A+\lambda B|^{1/n}
\ge(1-\lambda)|A|^{1/n}+\lambda|B|^{1/n}.
\]
If $K\subset\R^3$ is convex and its nonempty horizontal slices have finite area, prove that the square root of slice area is concave on their interval of heights. Distinguish this concavity conclusion from the displayed lower bound for arbitrary compact sets.
\end{exercise}

\begin{exercise}[$\star\star$]\label{ex:v9-app-j-long-needle}
Let a straight needle have length $\ell\ge0$ on unit-spaced lines.
(a) Prove that its crossing probability is the angular mean of $\min(1,\ell|\sin\theta|)$.
(b) Evaluate it for $\ell\le1$ and for $\ell>1$.
(c) Compare it with $\mathbb E[N]=2\ell/\pi$ and find the leading term of the missing probability as $\ell\to\infty$. Why can the expectation exceed $1$ without a contradiction?
\end{exercise}

\begin{exercise}[$\star\star$]\label{ex:v9-app-j-parallel-stability}
Let $K,M$ be nonempty compact convex planar sets with
$K\subset M+\delta\mathbb B$ and $M\subset K+\delta\mathbb B$, where $\delta\ge0$.
(a) Bound $|w_K(\theta)-w_M(\theta)|$ and $|L(K)-L(M)|$.
(b) Use Steiner's formula to bound $\bigl||K|-|M|\bigr|$ by an expression involving $\delta$ and the larger perimeter.
(c) Deduce continuity of both area and perimeter for such approximations, including polygons collapsing to a segment or point.
\end{exercise}

\begin{exercise}[$\star\star\star$]\label{ex:v9-app-j-smooth-width}
Put $u_\theta=(\cos\theta,\sin\theta)$, $v_\theta=(-\sin\theta,\cos\theta)$ and
$h(\theta)=1+\eta\cos3\theta$, where $0<|\eta|<1/8$.
(a) For $\gamma(\theta)=h(\theta)u_\theta+h'(\theta)v_\theta$, show $\gamma'=(h+h'')v_\theta$ and that the maximum of $\gamma(t)\cdot u_\theta$ over $t$ is uniquely attained at $t=\theta$ modulo $2\pi$.
(b) Prove that the convex hull of this curve has support function $h$ and constant width $2$, and that its boundary is the curve. Explain why it is not a disk.
(c) Compute its perimeter and area using the length integral and Green's theorem. Compare with the radius-$1$ disk.
\end{exercise}

\begin{exercise}[$\star\star$]\label{ex:v9-app-j-concavity-obstruction}
Let $A=[0,1]\cup[3,4]\subset\R$ and put $f(\lambda)=|(1-\lambda)A+\lambda A|$ for $0\le\lambda\le1$.
(a) Write the sum as four intervals and evaluate $f(0),f(1/3),f(1/2),f(2/3),f(1)$.
(b) Use these values to disprove concavity of $f$.
(c) Check that Brunn--Minkowski still gives a valid lower bound, and identify the convexity hypothesis that fails.
\end{exercise}

\section{Hints}

\begin{hint}{geometry:ex:helly1d}
Take the largest lower endpoint and smallest upper endpoint in the extended real line. If they coincide, use a pair attaining them to settle whether that endpoint is included.
\end{hint}
\begin{hint}{geometry:ex:hellysharp}
Take the facets of an $n$-simplex. Barycentric coordinates cannot all be zero.
\end{hint}
\begin{hint}{geometry:ex:hellydisk}
Regard a disk centre as the unknown point. The allowed centres for a point $p$ form the radius-$1$ disk centred at $p$.
\end{hint}
\begin{hint}{geometry:ex:caratheodory}
Start with a finite combination having the fewest positive coefficients. Use an affine dependence to subtract a multiple of its coefficients until one weight becomes zero.
\end{hint}
\begin{hint}{geometry:ex:lazzarini}
Reduce $C/T=2(5/6)/(355/113)$ to lowest terms. An arithmetic compatibility is not a probability-zero impossibility.
\end{hint}
\begin{hint}{geometry:ex:noodle}
The square's vertical extent is $a(|\sin\theta|+|\cos\theta|)$. A circle's boundary crosses a horizontal line twice when that height lies strictly between its top and bottom.
\end{hint}
\begin{hint}{geometry:ex:reuleauxarea}
Add three $60^\circ$ sectors and correct for the equilateral triangles counted more than once.
\end{hint}
\begin{hint}{geometry:ex:minkowskiwidth}
A linear projection takes a sum to the sum of projections. The sum of intervals has length equal to the sum of their lengths.
\end{hint}
\begin{hint}{geometry:ex:steiner}
Decompose the polygon's parallel region into edge rectangles and vertex sectors. For the limit, use nested dense inscribed polygons and prove that a compact convex set with interior has boundary of area zero by scaling it about an interior point.
\end{hint}
\begin{hint}{geometry:ex:constwidthdisk}
Substitute $L=\pi w$ into $4\pi A\le L^2$. When $w=0$, the set has only one point.
\end{hint}
\begin{hint}{geometry:ex:bmone}
After separate translations of the sets, arrange $\max A=0=\min B$. Both translated sets then lie in their sum and meet only at $0$.
\end{hint}
\begin{hint}{geometry:ex:bmconcave}
Scale the two sets before applying Brunn--Minkowski. Identify each horizontal slice with a subset of $\R^2$ and check the inclusion of weighted sums of slices.
\end{hint}
\begin{hint}{ex:v9-app-j-long-needle}
Conditional on the direction, a vertical interval of height $h$ hits an integer with probability $\min(1,h)$. Split the angular integral at $\arcsin(1/\ell)$ when $\ell>1$.
\end{hint}
\begin{hint}{ex:v9-app-j-parallel-stability}
Projection enlarges an interval by $2\delta$ when a radius-$\delta$ disk is added. For area, apply the two inclusions and Steiner's formula separately.
\end{hint}
\begin{hint}{ex:v9-app-j-smooth-width}
Set $\rho=h+h''=1-8\eta\cos3\theta>0$. For a fixed normal $u_\theta$, differentiate $\gamma(t)\cdot u_\theta$. For area compute $\det(\gamma,\gamma')/2$ and integrate.
\end{hint}
\begin{hint}{ex:v9-app-j-concavity-obstruction}
The four intervals start at $0,3\lambda,3(1-\lambda),3$ and all have length $1$. Compare the middle value with those at $1/3$ and $2/3$.
\end{hint}

\section{Solutions}

\begin{solution}{geometry:ex:helly1d}
\idea The innermost two endpoints decide the intersection; their inclusion matters when they coincide.
\writeup Let $a$ be the largest lower endpoint and $b$ the smallest upper endpoint, allowing infinite endpoints. Pairwise intersection implies $a\le b$. If $a<b$, choose any real point strictly between them; it belongs to every interval. If $a=b$, this common value is finite because the intervals are nonempty subsets of $\R$. Choose intervals attaining these two endpoints, possible because the family is finite. Their intersection forces the common endpoint to belong to both. Any other interval meets the first, so its upper endpoint cannot exclude the value $a$ if it equals $a$; it also meets the second, so its lower endpoint cannot exclude $a$ if it equals $a$. Intervals with endpoints strictly outside already contain $a$. Thus every interval contains it. This covers singleton intervals and all endpoint conventions.
\end{solution}

\begin{solution}{geometry:ex:hellysharp}
\idea Each facet omits one vertex, while all the other facets contain that vertex.
\writeup Choose affinely independent vertices $v_0,\ldots,v_n$ and let $F_i$ be the convex hull of all vertices except $v_i$. Any $n$ facets omit one $F_i$ and contain $v_i$. If a point lay in every facet, its unique barycentric coordinates in the simplex would satisfy $t_i=0$ for every $i$, contradicting $\sum_i t_i=1$. For $n=1$ these facets are two different singleton points, which gives the same example. For $n=0$, take the single empty convex set in $\R^0$: the intersection of the empty subfamily is $\R^0$, while the full intersection is empty.
\end{solution}

\begin{solution}{geometry:ex:hellydisk}
\idea Exchange the roles of points to cover and possible centres.
\writeup A centre $c$ covers $p$ by a closed radius-$1$ disk precisely when $|c-p|\le1$. Let $D_p$ be the set of these possible centres, a closed convex disk. For at least three points the hypothesis says every three $D_p$ intersect, so planar Helly gives a centre in their full intersection. For one or two points, the strengthened hypothesis applies to $S$ itself; for $S=\varnothing$ any disk works. Under the original wording ``any three'', a set of two points more than $2$ apart satisfies the hypothesis vacuously but cannot fit in one radius-$1$ disk. The at-most-three formulation repairs exactly this omission.
\end{solution}

\begin{solution}{geometry:ex:caratheodory}
\idea Shift weights along an affine dependence until the representation loses a point.
\writeup Choose a representation $x=\sum_{i=1}^m t_ix_i$ with the smallest possible number of positive weights, all $t_i>0$. If $m\ge n+2$, linear dependence of $(x_i,1)$ gives $\sum\lambda_ix_i=0$, $\sum\lambda_i=0$, with some $\lambda_i>0$. Put $c=\min_{\lambda_i>0}t_i/\lambda_i>0$. The new coefficients $t_i-c\lambda_i$ are nonnegative, sum to $1$, and represent the same $x$, but at least one is zero. This contradicts minimality, proving $m\le n+1$. In dimension zero any nonempty set is the singleton of $\R^0$, so one point suffices. If $S$ is empty, its hull has no points and the assertion is vacuous.
\end{solution}

\begin{solution}{geometry:ex:lazzarini}
\idea Solve the claimed estimate for the required rational crossing frequency.
\writeup The estimate equals $2\ell T/C=(5/3)T/C$. To obtain $355/113$ requires $C/T=113/213$. Since these integers are coprime, $T$ must be a positive multiple of $213$. For $T=3408=16\cdot213$, one needs $C=16\cdot113=1808$, and indeed $5680/1808=355/113$. The smallest compatible throw count is $213$. These specially compatible counts explain the suspicion, but do not logically rule out the observation under independent throws. Statistical assessment also needs the sampling and stopping rules; Badger discusses those issues.
\end{solution}

\begin{solution}{geometry:ex:noodle}
\idea Sum segment contributions for the square; for the circle count two intersections per intermediate height.
\writeup The square's perimeter is $4a$, so its mean count is $8a/\pi$. For a circle with any fixed orientation, each horizontal line strictly between its bottom and top meets it twice. The interval of heights has length $2r$, and the interval lemma therefore gives mean count $2(2r)=4r$; tangencies form a null set. The square has vertical extent $a(|\sin\theta|+|\cos\theta|)$, whose minimum is $a$. If $a\ge1$, every such interval contains an integer, so the boundary meets a ruled line in every placement. If $a<1$, place an unrotated square strictly between two lines; small changes of position and direction still miss. Thus the exact condition is $a\ge1$.
\end{solution}

\begin{solution}{geometry:ex:reuleauxarea}
\idea The three circular caps are sectors with an equilateral triangle removed.
\writeup A unit equilateral triangle has area $\sqrt3/4$. Each cap outside one of its sides has area $\pi/6-\sqrt3/4$. The Reuleaux area is consequently
\[
\frac{\sqrt3}{4}+3\left(\frac\pi6-\frac{\sqrt3}{4}\right)
=\frac{\pi-\sqrt3}{2}\approx0.7048.
\]
The same-width disk has area $\pi/4\approx0.7854$. Scaling lengths by $w$ multiplies areas by $w^2$, giving $(\pi-\sqrt3)w^2/2$ for the Reuleaux triangle.
\end{solution}

\begin{solution}{geometry:ex:minkowskiwidth}
\idea Add the projection intervals, then integrate their lengths.
\writeup The orthogonal projection $p_\theta$ is linear, so $p_\theta(K+M)=p_\theta(K)+p_\theta(M)$. Compactness and convexity make these projections closed intervals, possibly points. Their endpoints add, proving $w_{K+M}(\theta)=w_K(\theta)+w_M(\theta)$. Integrate Cauchy's formula to obtain $L(K+M)=L(K)+L(M)$. A radius-$\varepsilon$ disk has width $2\varepsilon$, so
$L(K+\varepsilon\mathbb B)=L(K)+2\pi\varepsilon$. The formula includes $\varepsilon=0$ and the segment and point perimeter conventions.
\end{solution}

\begin{solution}{geometry:ex:steiner}
\idea First compute the exact pieces of a polygon's outer layer, then justify the area limit rather than differentiating an approximation.
\writeup A point outside a compact convex set has a unique nearest point: if two different points minimized squared distance, their midpoint would give a strictly smaller value by the identity for the squared norm. For a polygon, the nearest point lies either on the interior of an edge or at a vertex. The first case supplies a rectangle of area $\ell_i\varepsilon$ along edge $i$; the second supplies a radius-$\varepsilon$ sector whose angle is that vertex's exterior angle $\phi_i$. These pieces have disjoint interiors. Since $\sum_i\phi_i=2\pi$, their total area is
\[
|K+\varepsilon\mathbb B|=|K|+\varepsilon\sum_i\ell_i
+\frac{\varepsilon^2}{2}\sum_i\phi_i
=A+L\varepsilon+\pi\varepsilon^2.
\]
The exterior-angle sum follows by following the successive outward normals once around a convex polygon, or by triangulating its interior angles. At $\varepsilon=0$ the formula is immediate.

Here are the details for passage to a compact convex $K$ with interior. Such a set has boundary of area zero. Translate an interior point to $0$ and take a disk $r\mathbb B\subset K$, $r>0$. For $0<t<1$, every point of $(1-t)K$ has a radius-$tr$ disk in $K$, by convexity. Thus $(1-t)K\subset\operatorname{int}K$, and
$|\operatorname{int}K|\ge(1-t)^2|K|$. Letting $t\to0$ proves $|\partial K|=0$.

Use the nested dense polygons $P_j$ from the Cauchy proof. Their union is dense in $K$ and contains its interior. To see the latter claim, surround an interior point by a small triangle in $K$; approximate its three vertices closely enough by points of the dense union. The point remains inside their triangle, and one $P_j$ contains all three approximations. Therefore $\mathbf1_{P_j}\to\mathbf1_K$ almost everywhere and $|P_j|\to|K|$.

For fixed $\varepsilon>0$, the convex nested sets $P_j+\varepsilon\mathbb B$ have union dense in $K+\varepsilon\mathbb B$. The same triangle argument puts every interior point in this union. The boundary-zero argument applies to this compact convex parallel set, so continuity from below gives
$|P_j+\varepsilon\mathbb B|\to|K+\varepsilon\mathbb B|$. We already know $L(P_j)\to L(K)$. Passing to the limit in the polygon identity proves Steiner's formula at each fixed $\varepsilon$, with no interchange of a derivative and a limit.

Finally a segment of length $a$ has parallel area $2a\varepsilon+\pi\varepsilon^2$, from a rectangle and two semicircles; a point has parallel area $\pi\varepsilon^2$. These agree with $L=2a$ and $0$. For $K$ with interior the difference quotient is $L+\pi\varepsilon$, whose limit is $L$, as required for the corollary.
\end{solution}

\begin{solution}{geometry:ex:constwidthdisk}
\idea Fixed width fixes perimeter, so the isoperimetric bound becomes an area bound.
\writeup If $w=0$, every projection is a point, forcing $K$ to be a point and its area to vanish. If $w>0$, Barbier gives $L=\pi w$. The isoperimetric inequality yields
$A\le L^2/(4\pi)=\pi w^2/4$. The disk of diameter $w$ has precisely this area and width, so it is a maximizer. This argument does not by itself classify all maximizers.
\end{solution}

\begin{solution}{geometry:ex:bmone}
\idea Put one set to the left of zero and the other to the right, with both touching zero.
\writeup Translate $A,B$ separately so that $\max A=0=\min B$; compact nonemptiness ensures these extrema exist. Translations do not change any of the three relevant measures. Since $0$ lies in both sets, $A\cup B\subset A+B$. Also $A\cap B=\{0\}$ has measure zero. Hence $|A+B|\ge|A\cup B|=|A|+|B|$, even when one or both measures vanish.
\end{solution}

\begin{solution}{geometry:ex:bmconcave}
\idea Convexity compares a weighted sum with a slice, while Brunn--Minkowski measures that sum.
\writeup Convexity gives $(1-\lambda)A+\lambda A\subset A$; writing $a=(1-\lambda)a+\lambda a$ proves the reverse inclusion, including $\lambda=0,1$. For nonempty compact $A,B$, apply Brunn--Minkowski to their scaled sets. The scaling law for volume gives the displayed lower bound; at the endpoints the zero-scaled set is $\{0\}$, so the same formula remains valid.

For the slice $K_t=\{(x,y):(x,y,t)\in K\}$, convexity yields
\[
(1-\lambda)K_s+\lambda K_t\subset K_{(1-\lambda)s+\lambda t}.
\]
The nonempty slices are convex and indexed by an interval. If the slices are compact, monotonicity of area and the planar Brunn--Minkowski inequality give the claim directly. In general convex sets are Lebesgue measurable, as explained in Appendix~\ref{geometry:ch:lattices}; take nonempty compact subsets of $K_s$ and $K_t$ whose areas approach their finite slice areas, using inner regularity from Volume~7. Their weighted sum is contained in the middle slice. Apply Brunn--Minkowski to those compact subsets and pass to the limit. Both cases give
\[
\sqrt{|K_{(1-\lambda)s+\lambda t}|}
\ge(1-\lambda)\sqrt{|K_s|}+\lambda\sqrt{|K_t|}.
\]
This also allows endpoint slices of area zero. It is concavity because $s,t$ are arbitrary. For arbitrary compact $A,B$, a lower bound by the chord joining the endpoint volumes alone is not a concavity theorem; Exercise~\ref{ex:v9-app-j-concavity-obstruction} gives a counterexample.
\end{solution}

\begin{solution}{ex:v9-app-j-long-needle}
\idea Probability counts whether an interval hits at all, while expectation counts how many integers it hits.
\writeup (a) At a fixed angle, the vertical extent is $h=\ell|\sin\theta|$. If $h\le1$, the set of offsets modulo $1$ causing a hit has measure $h$; if $h\ge1$, every interval of that length contains an integer. Endpoint exceptions have probability zero. Thus the conditional hit probability is $\min(1,h)$.
(b) For $0\le\ell\le1$, the average is $2\ell/\pi$. For $\ell>1$, put $a=\arcsin(1/\ell)$ and use quadrant symmetry:
\[
\Pr(N\ge1)=\frac2\pi\left(\int_0^a\ell\sin\theta\,d\theta
+\int_a^{\pi/2}1\,d\theta\right)
=\frac2\pi\left(\ell-\sqrt{\ell^2-1}+\arccos(1/\ell)\right).
\]
At $\ell=1$ this agrees with $2/\pi$.
(c) Since $N\ge\mathbf1_{N\ge1}$, expectation is at least the probability, and the segment calculation gives $\mathbb E[N]=2\ell/\pi$ for every length. For $\ell>1$, some directions and positions have multiple crossings, explaining the strict difference. Expanding $\sqrt{\ell^2-1}=\ell-1/(2\ell)+O(\ell^{-3})$ and $\arccos(1/\ell)=\pi/2-1/\ell+O(\ell^{-3})$ gives
$\Pr(N=0)=1/(\pi\ell)+O(\ell^{-3})$. An expected count can exceed $1$ because the count is not a zero--one variable.
\end{solution}

\begin{solution}{ex:v9-app-j-parallel-stability}
\idea Bound each quantity in both directions using the two parallel-set inclusions.
\writeup (a) Projection of the first inclusion gives $w_K\le w_M+2\delta$; the second gives the reverse bound. Hence $|w_K-w_M|\le2\delta$ pointwise. Cauchy's formula then gives $|L(K)-L(M)|\le2\pi\delta$.
(b) By inclusion and Steiner's formula,
\[
|K|-|M|\le L(M)\delta+\pi\delta^2,
\qquad |M|-|K|\le L(K)\delta+\pi\delta^2.
\]
Thus $\bigl||K|-|M|\bigr|\le\max(L(K),L(M))\delta+\pi\delta^2$.
(c) If one set is fixed and $\delta\to0$, part (a) bounds the other perimeters and makes their difference tend to zero. Part (b) then makes the area difference tend to zero. The formulas apply unchanged to a limiting segment or point, with areas zero and perimeters twice the segment length or zero.
\end{solution}

\begin{solution}{ex:v9-app-j-smooth-width}
\idea Prescribe each support line by its normal; positive curvature radius keeps the resulting curve convex.
\writeup (a) Since $u'=v$ and $v'=-u$, differentiation gives $\gamma'=(h+h'')v=\rho v$, where $\rho=1-8\eta\cos3\theta>0$. For a fixed normal $u_\theta$,
\[
\frac d{dt}\bigl(\gamma(t)\cdot u_\theta\bigr)=\rho(t)\sin(\theta-t).
\]
It is positive for $\theta-\pi<t<\theta$ and negative for $\theta<t<\theta+\pi$. Thus the unique maximum modulo $2\pi$ is at $t=\theta$, and its value is $h(\theta)$.
(b) The convex hull $K$ is compact, and a linear functional has the same maximum on a set and its convex hull. Hence its support function is $h$. The maximizing point for each normal is uniquely $\gamma(\theta)$. Every boundary point of a compact convex set with interior has a support line: for completeness, take outside points approaching it, project each to a nearest point, use the inequality $(x-p)\cdot(q-p)\le0$ obtained by differentiating squared distance along $[p,q]$, and take a convergent subsequence of the unit outward normals. Here $h>0$ and the curve has nonparallel tangents, so $K$ has interior. Every boundary point is therefore one of the unique maximizers, and each $\gamma(\theta)$ is on the boundary. The uniqueness from (a) also makes this regular closed parametrization injective modulo $2\pi$.

The width is $h(\theta)+h(\theta+\pi)=2$. It is not a disk: a disk with this width has radius $1$ and support function $1+c_1\cos\theta+c_2\sin\theta$, while the nonzero third harmonic $\eta\cos3\theta$ cannot have that form.
(c) Its speed is $\rho$, so $L=\int_0^{2\pi}\rho\,d\theta=2\pi$. Green's theorem, already available in Chapter~3, gives
\[
A=\frac12\int_0^{2\pi}\det(\gamma,\gamma')\,d\theta
=\frac12\int_0^{2\pi}h\rho\,d\theta
=\pi-4\pi\eta^2<\pi.
\]
Here the linear cosine terms integrate to zero and $\int_0^{2\pi}\cos^23\theta\,d\theta=\pi$. The disk of the same width has perimeter $2\pi$ and larger area $\pi$.
\end{solution}

\begin{solution}{ex:v9-app-j-concavity-obstruction}
\idea A weighted sum can develop and close gaps as the weight changes.
\writeup (a) Adding the four pairs of component intervals gives
\[
(1-\lambda)A+\lambda A
=[0,1]\cup[3\lambda,1+3\lambda]
\cup[3(1-\lambda),1+3(1-\lambda)]\cup[3,4].
\]
At $\lambda=0,1$ the union has measure $2$. At $1/3$ and $2/3$ the intervals fill $[0,4]$, giving measure $4$. At $1/2$ the middle two coincide in $[3/2,5/2]$, disjoint from the two outer intervals, so the measure is $3$.
(b) Concavity would require $f(1/2)\ge(f(1/3)+f(2/3))/2=4$, contradicting $3$.
(c) Brunn--Minkowski in dimension one still gives $f(\lambda)\ge(1-\lambda)|A|+\lambda|A|=2$. All the computed values satisfy it. The failed step is $(1-\lambda)A+\lambda A=A$: the set $A$ is not convex. Thus this example does not contradict the slice concavity for a convex body.
\end{solution}

\section*{Sources and further reading}

\begin{itemize}
\item Virgil, \emph{Aeneid}, Book~I, 365--368 (Dido and the bull's hide).
\item G.-L.~Leclerc de Buffon, \emph{Essai d'arithm\'etique morale}, in \emph{Histoire naturelle}, Suppl\'ement~4, 1777; his 1733 memoir in the \emph{Histoire
de l'Acad\'emie Royale des Sciences}; J.-\'E.~Barbier, \emph{Note sur le probl\`eme de l'aiguille et le jeu du joint couvert}, J.\ Math.\ Pures Appl.\ 5 (1860).
\item L.~Badger, \emph{Lazzarini's lucky approximation of $\pi$}, Mathematics Magazine 67 (1994), 83--91.
\item L.~Danzer, B.~Gr\"unbaum and V.~Klee, \emph{Helly's theorem and its relatives}, in \emph{Convexity}, Proc.\ Symp.\ Pure Math.\ 7, AMS, 1963 (history of
Helly's and Radon's theorems).
\item R.~J.~Gardner, \emph{The Brunn--Minkowski inequality}, Bulletin of the AMS 39 (2002), 355--405.
\item V.~Bl\aa sj\"o, \emph{The isoperimetric problem}, American Mathematical Monthly 112 (2005), 526--566 (Dido, Steiner, Weierstrass and Hurwitz).
\item D.~A.~Klain and G.-C.~Rota, \emph{Introduction to Geometric Probability}, Cambridge University Press, 1997 (Buffon, Cauchy and Crofton formulas).
\end{itemize}
